\subsection{A group-theoretic extension test in weight ten}
\label{subsec:eta-ten-extension-threshold}

\begin{proposition}\label{prop:eta-ten-no-extension-d7}
The character $\eta_{10}$ on $D_9$ extends invariantly to $D_8$
by $\alpha_{10}$, but has no $G$-invariant extension to $D_7$.
Consequently it has no such extension to $D_5$, $D_3$, $D$ or $N$.
\end{proposition}

\begin{proof}
The extension to $D_8$ was proved above. We give an obstruction
using only class-two commutator identities and the previously
established leading terms.

For $r\ge3$, let $\mathcal F^r$ be the subgroup of
$M=\ker\lambda$ whose pro-$2$ Magnus parameter starts in degree
at least $r$. The filtered normalized group law gives
\[
 [\mathcal F^r,\mathcal F^s]\subseteq\mathcal F^{r+s}.
\]
One has $M=\mathcal F^3$ and
$[M,M]\subseteq D_7\subseteq\mathcal F^8$. Thus
$[M,[M,M]]\subseteq\mathcal F^{11}$, and the quotient
$M/\mathcal F^{11}$ has nilpotence class at most two.
Moreover,
\[
 [\mathcal F^5,\mathcal F^5]\subseteq\mathcal F^{11}.
\]
Indeed its possible degree-ten leading term is the Ihara bracket
of two multiples of $\psi_5$, and this is zero. Vanishing over
$\mathbf Q_2$ implies vanishing in the integral leading module,
which is torsion-free.

Keep the arithmetic $\sigma,h$ of the weight-ten construction, and
put
\[
 a_0=\sigma,\qquad a_1=h\sigma h^{-1},\qquad
 a_2=h^2\sigma h^{-2},\qquad c_{ij}=[a_i,a_j].
\]
The weight-three law shows that
\[
 r=a_1a_0^{-27},\qquad q=a_2a_0^{-729}
\]
belong to $\mathcal F^5$. Therefore $[r,q]=1$ in
$M/\mathcal F^{11}$. All commutators are central in this quotient,
and expanding this identity gives
\[
 1=[a_1a_0^{-27},a_2a_0^{-729}]
   =c_{12}c_{01}^{729}c_{02}^{-27}.
\]
Hence $c_{12}=c_{02}^{27}c_{01}^{-729}$ there.

The arithmetic weight-eight and weight-ten witnesses are
\[
 w_8=c_{01},\qquad
 w_{10}=[a_0,(a_2a_1^{-27})(a_1a_0^{-27})^{-243}].
\]
The same class-two calculation gives
\[
 w_{10}=c_{02}c_{01}^{-270}\pmod{\mathcal F^{11}}.
\]
Consequently the word
\[
 z=h w_8h^{-1}w_8^{-6561}=c_{12}c_{01}^{-6561}
\]
satisfies
\[
 z=c_{02}^{27}c_{01}^{-7290}
   =w_{10}^{27}\pmod{\mathcal F^{11}}.
\]
Since $z\in J_7$ and its parameter agrees through degree ten with
that of $w_{10}^{27}$, one has $z\in J_9$ and
\[
 m_{10}(z)=27m_{10}(w_{10}),\qquad\eta_{10}(z)=1.
\]
On the other hand, every $G$-invariant character
$\chi:D_7\to\mathbf F_2$ satisfies
\[
 \chi(z)=\chi(h w_8h^{-1})-6561\chi(w_8)=0,
\]
because $6561$ is odd. Thus such a character cannot restrict to
$\eta_{10}$ on $D_9$. Any extension to one of the larger listed
subgroups would restrict to an extension on $D_7$, so those are
also impossible.
\end{proof}

\begin{corollary}\label{cor:weight-ten-character-kernels}
The invariant quadratic-character kernels on restriction to
$D_{10}$ are
\[
\begin{array}{c|l}
 \text{domain}&\text{kernel}\\ \hline
 D_9&\mathbf F_2\eta_{10}\\
 D_8&\mathbf F_2\eta_9\oplus\mathbf F_2\alpha_{10}\\
 D_7&\mathbf F_2\eta_8\oplus\mathbf F_2\eta_9\\
 D_5&\mathbf F_2\eta_7\oplus\mathbf F_2\alpha_8\\
 D_3&\mathbf F_2\eta_5\\
 N&\displaystyle\left(\bigoplus_{q\text{ odd prime}}\mathbf F_2\psi_q\right)
                    \oplus\mathbf F_2\eta_3.
\end{array}
\]
\end{corollary}

\begin{proof}
The first two rows follow from the exact scalar and joint maps.
For any larger listed domain, a character killed on $D_{10}$
restricts on $D_9$ to zero or $\eta_{10}$. The latter is impossible
by the preceding proposition. Such a character therefore already
kills $D_9$, and the previously proved weight-nine kernel table
applies.
\end{proof}
